% !TeX root = ../main.tex
% Add the above to each chapter to make compiling the PDF easier in some editors.

\chapter{Solution Approach}\label{chapter:solution-approach}

This chapter presents the technical approach developed to improve camera pose and intrinsic estimation by combining AnyCam's multi-view pose prediction with AnyCalib's single-view calibration capabilities. The core idea is to leverage multi-frame temporal consistency to produce more accurate and stable predictions than either method achieves alone.

\section{Overview and Design Philosophy}\label{sec:solution-overview}

The motivation for this work stems from a key observation about AnyCam: it distills depth and flow information into a pose prediction model, but as a byproduct, it must predict camera calibration. AnyCam handles this through an expensive 32-candidate focal length selection procedure, where all candidates are tested via flow reprojection and the best one is selected. This approach is computationally costly and, as shown in preliminary benchmarking experiments, can produce inaccurate calibration estimates.

AnyCalib, on the other hand, focuses exclusively on calibration. It predicts camera intrinsics from a single image by regressing per-pixel 3D ray directions and fitting a camera model in closed form. While more accurate than AnyCam's candidate system for single-frame predictions, AnyCalib ignores temporal information entirely---each frame is processed independently with no consistency enforcement across a video sequence.

This raises a central research question: \textbf{Can we combine both models to distill calibration, depth, and flow into a unified system that produces more accurate pose predictions and significantly better calibration than either model achieves independently?}

\subsection{Guiding Design Principles}

The solution approach is guided by several key principles:

\begin{itemize}
    \item \textbf{Adaptation over retraining}: As researchers working with large pretrained models, the goal is to make lightweight adapters or targeted modifications rather than retraining entire models from scratch. This approach is more efficient and aligns with best practices in modern deep learning research. All backbone models (DINOv2, UniDepth, UniMatch) remain frozen throughout training.
    % not true at all, we start with the pretrained models, and we do freeze them at first for initial phases of training (depends what we call each phase), but then during the end to end training, everything is trained end to end and unfrozen, including all backbones.

    \item \textbf{Staged training}: Begin with simple configurations and incrementally add complexity. Each stage is validated before proceeding to the next, ensuring that architectural choices are sound and that training converges reliably.
    
    \item \textbf{Self-supervised training}: The system must train without ground-truth labels for camera pose or intrinsics, using only the geometric constraints provided by flow reprojection. This matches AnyCam's original training paradigm and enables scaling to large, unlabeled video datasets.
    % however, for some phases of training, we generate our own so called pseudo ground truths. for example as we will see later, phase B involves training the new calibration FAT head, and during its phase 1, we teach the new FAT head to learn the average of anycalib calibration, by running anycalib over the sequence in question and using that as the initial pseudo labels, before then moving on to further training completely self supervised.

    \item \textbf{Multi-frame consistency}: Exploit the fact that a camera's intrinsic parameters are fixed across a video sequence. By processing multiple frames jointly and enforcing consistency, both calibration stability and pose accuracy can be improved.
\end{itemize}

\subsection{Overall Training Pipeline}

The complete training pipeline consists of three major phases:

\begin{enumerate}
    \item \textbf{Phase A: Pose Head Pre-training} \\
    Train a new pose head on AnyCam using static AnyCalib calibration (averaged per-frame predictions). All other AnyCam components remain frozen. This establishes a pose prediction baseline that uses AnyCalib's focal length instead of the 32-candidate system.

    \item \textbf{Phase B: Calibration Head Training} \\
    Train a dedicated calibration head (FAT architecture, described in detail later) to learn multi-frame calibration aggregation. This happens completely independently to the anycam pipeline and is unrelated to phase A completely. This is where we teach the calibration head to leverage multi-frame data in a more complex model architecture that will later allow for deeper understanding of the sequence as a whole. the satisfactory goal here is just maintain anycalib's accuracy, (improving upon it at this stage is a nice to have), while augmenting the complexity and depth of the model architecture, later enabling deeper learning over more features, such as pose information but especially multi frame info.

    \item \textbf{Phase C: End-to-End Joint Training} \\
    Connect the two pipelines (anycam with new pose head, and new fAT anycalib) into one another (specifically, we hook up the clabration prediction from fat calibration head into the projection matrix of anycam, which then predicts a pose. everything is unfrozen, including the backbones, with the exception of the depth and flow estimators unidepth and unimatch, but train in alternation, so one epoch the entire calibration pipeline is frozen while the anycam trains unfrozen end to end, and next epoch the other way around, anycam completely frozen including pose head, and allow anycalib fat to train end to end.
\end{enumerate}

The calibration head architecture developed for this work is \textbf{FAT} (Feature Aggregation Transformer), which aggregates multi-frame calibration information at the level of spatial features rather than final scalar outputs.

\section{AnyCalib Integration into AnyCam}\label{sec:anycalib-integration}

The first step in the solution approach is to replace AnyCam's expensive 32-candidate focal length system with direct predictions from AnyCalib.

\subsection{Architectural Modification}

In the original AnyCam architecture, the \texttt{sequence\_info\_head} predicts 32 discrete focal length candidates spanning a range from 0.1 to 4.0 times a reference focal length. During training, all 32 candidates are tested: for each candidate, the predicted poses are used to compute induced optical flow (via depth and camera motion), and the candidate that minimizes the flow reprojection error is selected.

This procedure is replaced with a single direct prediction from AnyCalib. AnyCalib takes an input image and predicts the full intrinsic matrix $(f_x, f_y, c_x, c_y)$ by:
\begin{enumerate}
    \item Extracting features using a pretrained DINOv2 Vision Transformer
    \item Predicting per-pixel 3D ray directions via a Dense Prediction Transformer (DPT) decoder
    \item Fitting a camera model to the predicted rays in closed form
\end{enumerate}

The result is a direct, model-agnostic calibration estimate that does not require candidate testing.

\subsection{Implementation via Wrapper Class}

The integration is implemented through an \texttt{AnyCamWrapperWithAnyCaLib} class that wraps the original AnyCam model. During the forward pass:
\begin{enumerate}
    \item AnyCalib runs on the input frame(s) to produce focal length predictions
    \item These predictions are injected into AnyCam's calibration matrix $\mathbf{K}$
    \item The normal AnyCam pose prediction pipeline proceeds using this fixed calibration
\end{enumerate}

For multi-frame sequences, AnyCalib produces per-frame predictions. The simplest baseline is to average these predictions to obtain a single calibration estimate per sequence. The calibration head architectures described in later sections aim to improve upon this naive averaging by learning to aggregate predictions more intelligently.

\subsection{Motivation from Benchmarking}

Preliminary benchmarking experiments (conducted in the \texttt{anycam-anycalib-benchmark/} directory) compared AnyCalib's focal length predictions against AnyCam's 32-candidate system on multiple datasets. The results showed that AnyCalib provides more stable and accurate predictions, with lower variance across frames and better alignment with ground truth when available. This empirical evidence motivated the integration approach and suggested that replacing the candidate system would improve overall system performance.

\section{Preliminary Experiments: Validating the Integration}\label{sec:preliminary-experiments}

Before developing full calibration head architectures, two preliminary experiments were conducted to validate the AnyCalib integration and explore the potential benefits of multi-frame consistency. These experiments serve as stepping stones that informed the final architectural design.

\subsection{Experiment 1: Single Pose Head with AnyCalib Focal Length}\label{sec:experiment-1}

\subsubsection{Objective}

The goal of Experiment 1 is to demonstrate that a pose head can learn to predict accurate camera poses when given AnyCalib's focal length predictions, thereby validating that AnyCalib is a viable replacement for AnyCam's 32-candidate system.

\subsubsection{Experimental Setup}

The experiment begins by loading the pretrained AnyCam model. The existing pose head is deleted and replaced with a freshly initialized head using random weights. All other components of AnyCam---including the transformer backbone, depth predictor, and flow estimator---are frozen. This results in approximately 16,897 trainable parameters, representing roughly 0.02\% of the total model parameters.

The frozen components ensure that learning is isolated to the pose head itself. AnyCalib's focal length prediction (from the first frame of each sequence) is used as the calibration input. Training is performed on the Objectron dataset, which contains 100 sequences of indoor object-centric videos. The dataset is split into train/validation/test sets with a 70/15/15 ratio, and the training set is used to extract all consecutive frame pairs.

The training objective is AnyCam's standard flow reprojection loss, which is fully self-supervised and requires no ground truth pose or calibration labels. The loss compares predicted optical flow (computed from predicted poses and depths) against observed optical flow from the frozen UniMatch model.

\subsubsection{Results}

The pose head successfully learned to predict poses. Training loss decreased from an initial value of 0.004339 to a final value of 0.001642, representing a 62\% improvement over the course of 50 epochs. This convergence demonstrates that the flow reprojection loss provides sufficient signal for the pose head to learn despite having random initialization.

To evaluate generalization, the trained model was tested on the Lightspeed validation dataset (200 samples with ground truth poses). The rotation error was 1.39 degrees mean and 1.01 degrees median, comparable to the AnyCam baseline (1.25 degrees mean, 0.91 degrees median). The slightly worse rotation accuracy is expected because only the pose head was retrained, whereas the calibration is now fixed rather than jointly optimized as in the original AnyCam training.

\subsubsection{Significance}

Experiment 1 confirms two critical hypotheses:
\begin{enumerate}
    \item AnyCalib provides a viable replacement for the 32-candidate focal length system
    \item A small, trainable pose head can achieve competitive accuracy using a frozen calibration source
\end{enumerate}

These findings validate the integration approach and establish a baseline for further improvements through multi-frame consistency.

\subsection{Experiment 2: Multi-Frame Consistency with Composed Flow}\label{sec:experiment-2}

\subsubsection{Objective}

Experiment 2 tests whether leveraging multiple frames and enforcing consistency across longer sequences can improve pose accuracy beyond the single-pair baseline established in Experiment 1.

\subsubsection{Flow Composition as a Key Innovation}

A central challenge in multi-frame training is computational efficiency. Computing optical flow using UniMatch is expensive, and running it on all pairs in a sequence (e.g., for 4 frames: 1$\to$2, 1$\to$3, 1$\to$4, 2$\to$3, 2$\to$4, 3$\to$4) would require $O(N^2)$ flow computations per sequence.

To address this, \textbf{flow composition} is introduced. Instead of running UniMatch on distant frame pairs, flows are composed through intermediate frames using bilinear interpolation. For example, the flow from frame 1 to frame 3 is computed by warping the flow from frame 2 to frame 3 through the flow from frame 1 to frame 2:
\begin{equation}
    \mathbf{w}_{1 \to 3}(u,v) = \mathbf{w}_{1 \to 2}(u,v) + \mathbf{w}_{2 \to 3}(\mathbf{p} + \mathbf{w}_{1 \to 2}(u,v))
\end{equation}
where $\mathbf{p} = (u,v)$ and the second term uses bilinear interpolation to sample the flow field at the warped location.

Flow composition is efficient (requires only $O(N)$ UniMatch computations for consecutive pairs) and maintains consistency with the consecutive flows. This approach enables multi-frame training without excessive computational overhead.

\subsubsection{Experimental Setup}

The setup is similar to Experiment 1, but each training sample now consists of $\text{max\_ahead} + 1$ frames. For example, with $\text{max\_ahead} = 3$, four frames are loaded per sequence. The model predicts poses for consecutive pairs (1$\to$2, 2$\to$3, 3$\to$4) and composes long-range poses mathematically:
\begin{equation}
    \mathbf{T}_{1 \to 3} = \mathbf{T}_{1 \to 2} \circ \mathbf{T}_{2 \to 3}
\end{equation}
where $\circ$ denotes pose composition in SE(3).

The flow reprojection loss is computed for both consecutive pairs (with full weight) and composed pairs (with 0.1$\times$ weight to balance loss magnitudes). This weighting ensures that the composed pairs contribute to learning without dominating the gradient signal.

A hyperparameter sweep was conducted, testing $\text{max\_ahead}$ values from 2 to 7. The results showed that $\text{max\_ahead} = 3$ (four frames total) is optimal. Shorter sequences ($\text{max\_ahead} = 2$) underutilize multi-frame information, while longer sequences ($\text{max\_ahead} \geq 4$) suffer from flow error accumulation during composition, which degrades performance.

\subsubsection{Results}

The best model (with $\text{max\_ahead} = 3$) achieved a rotation error of 1.20 degrees mean and 0.89 degrees median on the Lightspeed validation set. This outperforms both Experiment 1 (1.39/1.01 degrees) and the AnyCam baseline (1.25/0.91 degrees), confirming that multi-frame consistency improves pose accuracy.

The improvement is particularly notable in the median error, which decreased by 12\% compared to Experiment 1 and 2\% compared to the baseline. This suggests that multi-frame consistency benefits a large fraction of test samples, not just edge cases.

\subsubsection{Significance}

Experiment 2 demonstrates that:
\begin{enumerate}
    \item Multi-frame consistency improves pose accuracy over the single-pair baseline
    \item Flow composition is an effective and efficient alternative to running UniMatch on all frame pairs
    \item There is an optimal sequence length ($\text{max\_ahead} = 3$) that balances information gain against error accumulation
\end{enumerate}

These findings motivated the development of dedicated calibration head architectures that can learn to aggregate multi-frame calibration information more intelligently than simple averaging. If multi-frame information improves pose predictions, it should also improve calibration predictions when leveraged appropriately.

\section{Calibration Head Architecture: FAT (Feature Aggregation Transformer)}\label{sec:fat-architecture}

The Feature Aggregation Transformer (FAT) represents a more sophisticated approach to multi-frame calibration, operating on rich spatial features rather than final scalar outputs.

\subsection{Design Rationale}

FAT is motivated by a key insight: AnyCalib's DINOv2 backbone is fine-tuned end-to-end during calibration training, so its intermediate features encode calibration-specific geometric information. By aggregating these \emph{spatial features} $[N, 1024, H/14, W/14]$ before they are collapsed to scalars, FAT preserves geometric relationships that would be lost in scalar averaging.

This can be understood through an analogy: scalar aggregation (as in DA3) is like averaging four people's measurements of a room (``the room is 5 meters wide on average''), while spatial aggregation is like combining four detailed 3D scans of the room and then extracting measurements from the combined scan. The latter preserves much more geometric structure and should lead to more accurate final estimates.

Importantly, FAT operates \emph{between} AnyCalib's Step 1 (DINOv2 feature extraction) and Step 2 (DPT decoder + ray head). This positioning allows FAT to aggregate calibration-tuned features while still producing the same output format (3D rays) that AnyCalib's calibrator expects.

\subsection{Pipeline Position}

FAT's position in the AnyCalib pipeline is:

\begin{enumerate}
    \item \textbf{Input}: $N$ frames $[\mathbf{I}_1, \ldots, \mathbf{I}_N]$ where $\mathbf{I}_i \in \mathbb{R}^{3 \times H \times W}$
    \item \textbf{Step 1: DINOv2 Backbone (FROZEN)}: Extract multi-scale features
    \begin{equation}
        \mathbf{F}_i^{(s)} \in \mathbb{R}^{1024 \times h_s \times w_s} \quad \text{for } s \in \{1,2,3,4\}
    \end{equation}
    where $s$ indexes scales from transformer blocks 4, 11, 17, 23.
    \item \textbf{Step 1.5: FAT (TRAINABLE)}: Aggregate across frames
    \begin{equation}
        \mathbf{F}^{(s)}_{\text{agg}} = \text{FAT}([\mathbf{F}_1^{(s)}, \ldots, \mathbf{F}_N^{(s)}]) \in \mathbb{R}^{1024 \times h_s \times w_s}
    \end{equation}
    \item \textbf{Step 2: DPT Decoder + Ray Head (FROZEN)}: Convert to 3D unit rays
    \begin{equation}
        \mathbf{R} = \text{RayHead}(\text{DPT}([\mathbf{F}^{(1)}_{\text{agg}}, \ldots, \mathbf{F}^{(4)}_{\text{agg}}])) \in \mathbb{R}^{H \times W \times 3}
    \end{equation}
    \item \textbf{Step 3: Calibrator (Non-differentiable)}: Fit intrinsics via RANSAC + least squares
    \begin{equation}
        \mathbf{c} = (f_x, f_y, c_x, c_y) = \text{Calibrator}(\mathbf{R})
    \end{equation}
\end{enumerate}

By inserting FAT between Steps 1 and 2, the model aggregates features while keeping the rest of AnyCalib's pipeline unchanged.

\subsection{FAT Module Architecture}

FAT processes each scale independently using the same aggregation mechanism. For a given scale with features $[\mathbf{F}_1, \ldots, \mathbf{F}_N]$ where $\mathbf{F}_i \in \mathbb{R}^{1024 \times h \times w}$:

\begin{enumerate}
    \item \textbf{Reshape}: Convert to spatial-first format
    \begin{equation}
        \mathbf{F}_{\text{reshaped}} \in \mathbb{R}^{(h \cdot w) \times N \times 1024}
    \end{equation}
    Each of the $h \cdot w$ spatial positions now has a sequence of $N$ frame tokens.

    \item \textbf{Self-Attention}: Apply 2 layers of multi-head self-attention (8 heads) across the $N$ frame tokens \emph{at each spatial position independently}. This allows each spatial location to exchange information across frames but maintains spatial independence.

    \item \textbf{Mean Pooling}: Aggregate across frames
    \begin{equation}
        \mathbf{F}_{\text{agg}}(p) = \frac{1}{N} \sum_{i=1}^N \mathbf{F}_i(p) \quad \text{for each position } p
    \end{equation}

    \item \textbf{Reshape Back}: Convert to original spatial format
    \begin{equation}
        \mathbf{F}_{\text{agg}} \in \mathbb{R}^{1024 \times h \times w}
    \end{equation}
\end{enumerate}

This design keeps computation manageable by operating on small sequences ($N=4$ frames per spatial position) rather than on the full spatial resolution ($h \times w$ positions).

\textbf{Optional Visual Conditioning}: In Phase 2 training, DINOv2-small CLS tokens $[\mathbf{v}_1, \ldots, \mathbf{v}_N]$ where $\mathbf{v}_i \in \mathbb{R}^{384}$ can be projected to 1024-dim and appended to each spatial position's frame tokens. This gives the transformer both per-position spatial features and global visual context. However, experiments showed that visual tokens caused overfitting when trained only on Objectron, and they were disabled for Phase 3.

\subsection{The Non-Differentiability Challenge}

A critical challenge in training FAT is that AnyCalib's calibrator uses RANSAC, which is non-differentiable. This means gradients cannot flow from the calibration output back through the calibrator to FAT's parameters. Three solutions were developed and tested:

\subsubsection{V1: Soft Exponential Weights (Failed)}

The first attempt was to make the loss differentiable by replacing RANSAC's hard inlier selection with soft exponential weighting. The approach:
\begin{enumerate}
    \item Compute predicted rays $\mathbf{R}^{\text{pred}}$ and expected rays $\mathbf{R}^{\text{exp}}$ from the current calibration estimate
    \item Compute angular error: $\theta(p) = \arccos(\mathbf{R}^{\text{pred}}(p) \cdot \mathbf{R}^{\text{exp}}(p))$
    \item Weight each ray by a smooth exponential function:
    \begin{equation}
        w(p) = \exp\left(-\frac{\theta(p)^2}{2\sigma^2}\right)
    \end{equation}
    \item Compute weighted loss: $\mathcal{L} = \sum_p w(p) \|\mathbf{R}^{\text{pred}}(p) - \mathbf{R}^{\text{exp}}(p)\|^2$
\end{enumerate}

This produced NaN gradients from batch 0 onward. The root cause was numerical instability in the composition of $\exp()$ and $\arccos()$, particularly in mixed-precision (FP16) training. The gradient path involves:
\begin{equation}
    \frac{\partial \mathcal{L}}{\partial \mathbf{R}} = \frac{\partial}{\partial \mathbf{R}} \left[ \exp\left(-\frac{\arccos(\cdots)^2}{2\sigma^2}\right) \right]
\end{equation}
which combines the numerically sensitive inverse-trigonometric derivative with the exponential function's tendency to produce very large or very small values. Even in full FP32 with gradient clipping, the gradient path remained unstable.

\subsubsection{V2: Implicit Differentiation Through Weighted Least Squares}

The second approach uses RANSAC for inlier selection (no gradients) but makes the subsequent least squares solve differentiable:

\begin{enumerate}
    \item Run RANSAC to obtain an inlier mask $\mathbf{m} \in \{0,1\}^{H \times W}$ (treated as a constant)
    \item Create weights: $w(p) = \begin{cases} 1.0 & \text{if } m(p) = 1 \\ 10^{-6} & \text{otherwise} \end{cases}$
    \item Formulate the calibration fitting as a weighted least squares problem:
    \begin{equation}
        \min_{\mathbf{x}} \|\mathbf{A}\mathbf{x} - \mathbf{b}\|_{\mathbf{W}}^2
    \end{equation}
    where $\mathbf{x} = (p, q, r, s)$ parameterizes intrinsics as $p = 1/f_x$, $q = c_x/f_x$, $r = 1/f_y$, $s = c_y/f_y$.
    \item Solve via: $\mathbf{x}^* = (\mathbf{A}^\top \mathbf{W}^2 \mathbf{A})^{-1} \mathbf{A}^\top \mathbf{W}^2 \mathbf{b}$
\end{enumerate}

The key insight is that $\mathbf{A}$ (pixel coordinates) and $\mathbf{W}$ (RANSAC-derived weights) are \emph{constants} with no gradients, while $\mathbf{b}$ (derived from rays) carries gradients. The linear algebra operations (matrix multiply, solve) are inherently stable and differentiable, so gradients flow:
\begin{equation}
    \mathcal{L} \to \mathbf{x}^* \to \mathbf{b} \to \mathbf{R} \to \text{FAT parameters}
\end{equation}

This approach works and is used in FAT Phase 3, but it requires careful implementation to ensure numerical stability in the matrix solve.

\subsubsection{V3: Direct Reprojection Loss (Used for Phases 1--2)}

The third and simplest approach avoids the calibrator entirely:

\begin{enumerate}
    \item Run AnyCalib on each frame independently to obtain per-frame predictions $[\mathbf{c}_1, \ldots, \mathbf{c}_N]$ (detached, no gradients)
    \item Compute average: $\bar{\mathbf{c}} = \frac{1}{N} \sum_i \mathbf{c}_i$ (detached)
    \item Use FAT to predict rays $\mathbf{R}^{\text{FAT}}$
    \item Project rays back to the image plane using $\bar{\mathbf{c}}$:
    \begin{align}
        u^{\text{proj}}(p) &= \bar{f}_x \cdot \frac{R_x(p)}{R_z(p)} + \bar{c}_x \\
        v^{\text{proj}}(p) &= \bar{f}_y \cdot \frac{R_y(p)}{R_z(p)} + \bar{c}_y
    \end{align}
    \item Compute MSE against the actual pixel grid
\end{enumerate}

Gradients flow directly: $\mathcal{L} \to (u^{\text{proj}}, v^{\text{proj}}) \to \mathbf{R}^{\text{FAT}} \to \text{FAT parameters}$. This is simpler, faster, and more numerically stable than V2, and it was used successfully for FAT Phases 1 and 2.

\subsection{FAT Training Phases}

\subsubsection{Phase 1: Feature Aggregation Pre-training}

\textbf{Objective}: Train FAT in isolation to verify that it can aggregate multi-frame features and produce accurate calibration.

\textbf{Setup}: Train only FAT (approximately 25M parameters). All other components frozen (DINOv2-Large backbone, DPT decoder, ray head). No visual tokens.

\textbf{Loss}: V3 reprojection loss (described above).

\textbf{Hyperparameters}: Learning rate 5$\times 10^{-5}$, batch size 2, max\_ahead=3 (4 frames), 50 epochs. Data sampling uses step=2 (every other frame pair) to speed up training while maintaining video diversity.

\textbf{Result}: The model converged successfully with stable training. Validation loss decreased steadily, and per-epoch benchmarking showed improving calibration accuracy. This confirmed that FAT can aggregate spatial features effectively.

\subsubsection{Phase 2: Visual-Conditioned Aggregation (Abandoned)}

\textbf{Objective}: Add visual conditioning via DINOv2-small CLS tokens to see if global visual context improves calibration.

\textbf{Setup}: Load Phase 1 checkpoint and unfreeze the visual projection layer. Add DINOv2-small CLS tokens as inputs.

\textbf{Loss}: V3 reprojection loss (same as Phase 1).

\textbf{Result}: Validation loss immediately increased (8.61 $\to$ 14.37 $\to$ 19.10 across epochs) while training loss decreased, indicating clear overfitting. This is likely due to the small size and homogeneity of the Objectron dataset (only indoor, object-centric videos). Visual tokens provide additional capacity but not enough diverse data to learn generalizable visual-calibration relationships.

\textbf{Decision}: Phase 2 was abandoned. Phase 3 loads directly from Phase 1, skipping visual conditioning.

\subsubsection{Phase 3: Combined Loss End-to-End Training}

\textbf{Objective}: Train FAT and a fresh pose head jointly using both flow reprojection and calibration reprojection losses.

\textbf{Setup}: Load FAT from Phase 1. Initialize a new pose head with random weights. Both are trainable; all backbones remain frozen.

\textbf{Loss}: Combined loss with two terms:
\begin{equation}
    \mathcal{L}_{\text{total}} = \lambda_{\text{flow}} \mathcal{L}_{\text{flow}} + \lambda_{\text{calib}} \mathcal{L}_{\text{calib}}
\end{equation}
where:
\begin{itemize}
    \item $\mathcal{L}_{\text{flow}}$ (weight $\lambda_{\text{flow}} = 1.0$): AnyCam's flow reprojection loss, the primary signal for pose learning. FAT's predicted focal length enters through the calibration matrix $\mathbf{K}$.
    \item $\mathcal{L}_{\text{calib}}$ (weight $\lambda_{\text{calib}} = 10^{-4}$): Phase 1's V3 reprojection loss, acting as an anchor to prevent calibration divergence. Without this term, FAT could produce extreme intrinsics that minimize flow error but are geometrically unreasonable.
\end{itemize}

The dual gradient path provides stability:
\begin{align}
    \mathcal{L}_{\text{flow}} &\to \text{induced flow} \to \text{poses (Pose Head)} \to f_x \text{ (FAT)} \\
    \mathcal{L}_{\text{calib}} &\to \text{projected coords} \to \text{FAT rays} \to \text{FAT}
\end{align}

Both FAT and the pose head receive gradients from $\mathcal{L}_{\text{flow}}$. FAT additionally receives gradients from $\mathcal{L}_{\text{calib}}$ through the ray projection path. This ensures that calibration predictions remain geometrically plausible while the pose head learns.

\textbf{Training}: Uses mixed precision (FP16 for frozen backbones, FP32 for FAT) with gradient scaling. Training is joint (not alternating), allowing both components to co-adapt from the start.

\textbf{Hyperparameters}: Learning rate $10^{-5}$, batch size 1--2, max\_ahead=3, gradient clipping with max norm 1.0.

\subsection{Summary}

FAT represents a sophisticated approach to multi-frame calibration aggregation. By operating on spatial features, it preserves geometric structure that is lost in scalar aggregation. The three-version loss evolution (V1 failed, V2 designed but complex, V3 adopted) demonstrates the challenges of working with non-differentiable components and the importance of finding simple, stable solutions. Phase 3's combined loss design enables end-to-end training while maintaining calibration stability through the anchor term.

\section{FAT Design Summary}\label{sec:fat-summary}

FAT represents a sophisticated approach to multi-frame calibration aggregation by operating on spatial features rather than final scalar outputs. The key advantages of this design are:

\begin{itemize}
    \item \textbf{Rich geometric information}: By aggregating calibration-tuned spatial features $[N, 1024, h, w]$ before they collapse to scalars, FAT preserves geometric relationships that enable more accurate final predictions.

    \item \textbf{Direct differentiability}: Operating on rays provides a more direct gradient path from the loss to the aggregation module, compared to scalar approaches where focal length enters only indirectly through the calibration matrix.

    \item \textbf{Multi-scale processing}: FAT aggregates features at four scales from the DINOv2 backbone, capturing both fine-grained local structure and coarse global patterns.

    \item \textbf{Combined loss training}: Phase 3's dual loss design (flow reprojection + calibration anchor) enables end-to-end pose and calibration learning while preventing geometric divergence.
\end{itemize}

The evolution through three loss formulations (V1 failed, V2 designed, V3 adopted) demonstrates the importance of numerical stability and gradient flow in self-supervised geometric learning. The final combined loss in Phase 3 successfully balances pose accuracy and calibration consistency.

\section{Overall Training Pipeline for Cluster Deployment}\label{sec:overall-pipeline}

The complete training pipeline that will be executed on the cluster integrates all components described above:

\subsection{Phase A: Pose Head Pre-training}

Using the approach validated in Experiment 2 (Section~\ref{sec:experiment-2}), train a new pose head on AnyCam with static AnyCalib calibration (average of per-frame predictions). Freeze all AnyCam components except the pose head. Train on all four datasets: RealEstate10K, YouTube VOS, WalkingTours, and OpenDV. Use max\_ahead=3 with composed flow losses. This gives the pose head a good initialization for the multi-frame setting and establishes a baseline pose prediction capability.

\subsection{Phase B: Calibration Head Training}

Train the FAT calibration head through its three phases while freezing AnyCam and the pose head from Phase A. Phase 1 trains FAT in isolation using reprojection loss. Phase 2 (abandoned due to overfitting on small datasets) attempted to add visual tokens. Phase 3 trains FAT and pose head jointly with combined loss.

This staged approach ensures that the calibration head learns to produce good calibration without interference from pose head updates in early phases, before final joint optimization.

\subsection{Phase C: End-to-End Joint Training}

Unfreeze both the calibration head and pose head, training them together with the combined loss:
\begin{equation}
    \mathcal{L}_{\text{total}} = \mathcal{L}_{\text{flow}} + \lambda_{\text{calib}} \mathcal{L}_{\text{calib}}
\end{equation}

All backbone models (DINOv2, UniDepth, UniMatch) remain frozen throughout. The combined loss allows the calibration and pose components to co-adapt, with the calibration anchor preventing divergence.

\subsection{Training Datasets}

Training uses the same four datasets as AnyCam's original work:
\begin{itemize}
    \item \textbf{RealEstate10K}: Real estate walkthrough videos with diverse indoor environments
    \item \textbf{YouTube VOS}: General YouTube videos covering many scene types
    \item \textbf{WalkingTours}: Urban walking videos with dynamic street scenes
    \item \textbf{OpenDV}: Driving videos with forward-facing camera motion
\end{itemize}

This diversity ensures that the model learns to handle varied camera motion, scene content, and lighting conditions. The total number of sequences will be determined after preprocessing completes.

\section{Summary}

This chapter presented the solution approach for improving camera pose and intrinsic estimation through multi-frame consistency and AnyCalib integration. The key contributions are:

\begin{enumerate}
    \item \textbf{AnyCalib integration}: Replacing AnyCam's expensive 32-candidate focal length system with direct predictions from AnyCalib, validated through preliminary experiments showing competitive pose accuracy (1.39° rotation error).

    \item \textbf{Multi-frame consistency}: Leveraging temporal information via composed flow losses (max\_ahead=3 optimal), improving pose accuracy to 1.20° rotation error---outperforming both single-pair baseline and AnyCam.

    \item \textbf{FAT calibration head}: A feature aggregation architecture operating on calibration-tuned spatial features rather than scalars, with direct differentiability through rays and a combined loss design for stable end-to-end training.

    \item \textbf{Staged training pipeline}: Three-phase approach (pose pre-training, calibration head training, joint optimization) ensuring robust learning and enabling systematic evaluation of each component.
\end{enumerate}

\textbf{Note on diagrams}: Several sections of this chapter (FAT pipeline position, multi-scale feature aggregation, combined loss gradient flow) would benefit from visual diagrams to complement the mathematical descriptions. These will be added in the final thesis version.
